\Def \textit{Префиксная сумма} по $j$ элементам - это сумма первых $j$ элементов массива. 

\Task Есть массив из $n$ элементов $a_1, \ldots, a_n$ и $q$ запросов вида $(l_1, r_i), ..., (l_q, r_q)$. Требуется для каждого запроса вывести $\sum\limits_{j = l_i}^{r_i}a_j = a_{l_i} + a_{l_i + 1} + ... + a_{r_i}$, где $i$ - номер запроса. 

\Alg Посчитаем префиксные суммы для любого $j$ от $0$ до $n$: $pref_j = a_1 + ... + a_j$. Для $j = 0: pref_0 = 0$, для $j > 0: pref_j = pref_{j-1} + a_j$. Ответ на запрос $(l_i, r_i):$ $\sum\limits_{j = l_i}^{r_i}a_j = pref_{r_i} - pref_{l_i - 1}$.

\textit{Асимптотика:} Подсчёт всех значений $pref_j$ (по порядку) займёт $O(n)$ времени, подсчёт ответа на все запросы: $O(q)$. Итоговая асимптотика $O(n + q)$. 

\textbf{Обобщение на произвольную обратимую операцию}

Подобный прием можно использовать не только для сложения, но и для других операций, являющихся обратимыми:
\begin{enumerate}
    \item побитовое исключающее «или» $\oplus$
    \item сложение по модулю
    \item умножение
\end{enumerate}

\textbf{Алгоритм для произвольной операции.}

Посчитаем префиксные значения для любого $j$ от $0$ до $n$: 

$$pref_j = \text{op}(a_1, a_2, \ldots, a_j).$$

Для $j = 0$: $pref_0 = \text{identity}$, где $\text{identity}$ — это нейтральный элемент для операции $\text{op}$ (например, $0$ для сложения, $1$ для умножения). Для $j > 0$: 

$$pref_j = \text{op}(pref_{j-1}, a_j).$$

Ответ на запрос $(l_i, r_i)$: 

$$\text{op}\left(pref_{r_i}, \text{inverse}(pref_{l_i - 1})\right),$$

где $\text{inverse}(x)$ — это функция, возвращающая обратное значение для операции $\text{op}$ (например, $-x$ для сложения, $\frac{1}{x}$ для умножения).

\pagebreak